\section{Pregunta N$^{\circ}$23\qquad Carlos Daniel Malvaceda Canales}

\begin{frame}
	\begin{enumerate}\setcounter{enumi}{22}
		\item

		      Dos kilos de plátanos y tres de peras cuestan S/$8,80$.
		      Cinco kilos de plátanos y cuatro de peras cuestan
		      S/$16,40$.
		      Determine el costo de kilo del plátano y de la pera según
		      el siguiente requerimiento.
		      \begin{enumerate}[a)]
			      \item

			            Modele el problema.

			      \item

			            Determine la norma matricial de $A$ y $A^{-1}$.

			      \item

			            Determine el condicionamiento de $A$.

			      \item

			            Indique si está bien o mal condicionado.
		      \end{enumerate}
	\end{enumerate}
	\begin{solution}

		\begin{enumerate}[a)]
			\item

			      .

			\item

			      .

			\item

			      .
		\end{enumerate}
	\end{solution}
\end{frame}
